\chapter{Für wen es wirkt}

\section{Der Durchschnitt verbirgt}

\begin{merksatz}
\boxtitel{Ein Effekt, vier Zahlen}
$300\,000$ Patienten, Behandlung ausgelost. Der Effekt hängt vom Alter ab: bei den Jüngsten $-0{,}20$, bei den Ältesten $0$.

\begin{center}
{\small
\renewcommand{\arraystretch}{1.48}
\begin{tabular}{@{}P{2.2cm}P{2.4cm}P{2.4cm}@{}}
\toprule
\textbf{Alter} & \textbf{geschätzt} & \textbf{wahr}\\
\midrule
$40$--$55$ & $-0{,}171$ & $-0{,}170$\\
$55$--$70$ & $-0{,}107$ & $-0{,}110$\\
$70$--$85$ & $-0{,}053$ & $-0{,}050$\\
$85$--$90$ & $-0{,}029$ & $-0{,}010$\\
\midrule
\textbf{alle} & \textbf{$-0{,}100$} & $-0{,}100$\\
\bottomrule
\end{tabular}}
\end{center}

\begin{center}
\begin{minipage}{0.9\textwidth}
\textbf{Der Durchschnittseffekt von zehn Punkten gilt für niemanden: Die Jungen gewinnen doppelt so viel, die Alten nichts.}
\end{minipage}
\end{center}

\begin{center}
\begin{minipage}{0.88\textwidth}
\kursiv{\textbf{Das ist die zweite Warnung des Vorworts in Zahlen.} Ein Durchschnittseffekt ist ein Mittelwert über eine Bevölkerung --- und wie jeder Mittelwert sagt er nichts über die Streuung. Er kann sogar positiv sein, während die Behandlung der Hälfte schadet; das sieht man ihm nicht an.}
\end{minipage}
\end{center}
\end{merksatz}

\begin{merksatz}
\boxtitel{Was man schätzen will}
\[ \tau(\vek z) = \mathbb{E}\bigl[Y(1) - Y(0) \ \big|\ Z = \vek z\bigr] \]
--- der Effekt bei Patienten mit den Merkmalen $\vek z$.

\begin{center}
{\small
\renewcommand{\arraystretch}{1.52}
\begin{tabular}{@{}P{2.6cm}P{6.0cm}@{}}
\toprule
\textbf{Nicht} & der Effekt \kursiv{dieses} Patienten --- der bleibt nach Satz 2.2 unbekannt.\\
\textbf{Sondern} & der Durchschnitt über alle, die ihm in den gemessenen Merkmalen gleichen.\\
\bottomrule
\end{tabular}}
\end{center}

\begin{center}
\begin{minipage}{0.88\textwidth}
\kursiv{\textbf{Je mehr Merkmale, desto näher kommt man dem Einzelnen --- und desto größer wird die Streuung}, denn in jeder Zelle bleiben weniger Patienten. Es ist der Zielkonflikt aus Band 22, Kapitel~2, nun auf eine Differenz angewandt: Man schätzt nicht einen Ausgang, sondern den Unterschied zweier Ausgänge, und Differenzen sind erheblich unschärfer als ihre Bestandteile.}
\end{minipage}
\end{center}
\end{merksatz}

\section{Wie man es schätzt}

\begin{merksatz}
\boxtitel{Drei Wege}
\begin{center}
{\small
\renewcommand{\arraystretch}{1.55}
\begin{tabular}{@{}P{2.8cm}P{5.8cm}@{}}
\toprule
\textbf{Untergruppen} & Vorher festgelegte Gruppen getrennt auswerten. Ehrlich, aber grob und schnell erschöpft.\\
\textbf{Wechselwirkung} & Ein Term $X\cdot Z$ im Modell. Unterstellt eine Form.\\
\textbf{Kausale Wälder} & Bäume, die nicht nach dem Ausgang, sondern nach dem \kursiv{Unterschied} der Ausgänge teilen --- Band 22, Kapitel~9.\\
\bottomrule
\end{tabular}}
\end{center}

\begin{center}
\begin{minipage}{0.88\textwidth}
\kursiv{Ein kausaler Wald sucht die Aufteilung, bei der sich die geschätzten Effekte der beiden Gruppen am stärksten unterscheiden. \textbf{Der entscheidende Kunstgriff heißt Ehrlichkeit}: Eine Hälfte der Daten bestimmt, wo geteilt wird, die andere schätzt in den entstandenen Blättern. Ohne diese Trennung findet der Wald Unterschiede, die er selbst erzeugt hat --- das Leck aus Band 22, Kapitel~3, in kausaler Gestalt.}
\end{minipage}
\end{center}
\end{merksatz}

\section{Der Bruchfall}

\begin{warnung}
\boxtitelr{Untergruppen, die keine sind}
\begin{center}
{\small
\renewcommand{\arraystretch}{1.52}
\begin{tabular}{@{}P{2.8cm}P{5.8cm}@{}}
\toprule
\textbf{Multiples Testen} & Zwanzig Untergruppen ergeben bei reinem Zufall eine \enquote{bedeutsame}.\\
\textbf{Skalenabhängig} & Ein Effekt kann auf der Differenzskala schwanken und auf der Verhältnisskala konstant sein --- und umgekehrt.\\
\textbf{Streuung} & Der Effekt in einer Untergruppe wird aus einer Differenz zweier Differenzen geschätzt --- vierfache Fallzahl für gleiche Genauigkeit.\\
\bottomrule
\end{tabular}}
\end{center}

\schritt{Zur zweiten Zeile}
\begin{center}
\begin{minipage}{0.88\textwidth}
\kursiv{Senkt eine Behandlung das Risiko überall um $20$ Prozent \kursiv{relativ}, so beträgt die absolute Senkung bei einem Grundrisiko von $40$ Prozent acht Punkte, bei zehn Prozent nur zwei. \textbf{Auf der Verhältnisskala gibt es dann keine Heterogenität, auf der Differenzskala erhebliche} --- und beide Beschreibungen sind richtig. Für die Behandlungsentscheidung zählt die absolute, denn sie sagt, wie viele Patienten man behandeln muss.}
\end{minipage}
\end{center}

\begin{center}
\begin{minipage}{0.88\textwidth}
\kursiv{\textbf{Daraus folgt: Die wichtigste Quelle von Heterogenität ist das Grundrisiko}, nicht eine besondere Empfindlichkeit. Wer wissen will, wem eine Behandlung am meisten nützt, schätzt zuerst das Risiko ohne Behandlung --- das ist eine Vorhersageaufgabe im Sinne von Band 22 --- und wendet darauf den relativen Effekt an. Dieses zweistufige Vorgehen ist robuster als die unmittelbare Schätzung von $\tau(\vek z)$.}
\end{minipage}
\end{center}
\end{warnung}

\section{Übungsaufgaben}

\begin{aufgabe}
\boxtitel{Aufgabe 14.1 --- Wie viele Patienten je Untergruppe?}
Der Gesamteffekt beträgt $-0{,}10$ bei einem Grundrisiko von $0{,}30$. Wie viele Patienten braucht man, um einen Unterschied von $0{,}10$ zwischen zwei Untergruppen nachzuweisen?

\loesung
\schritt{Schritt 1 --- Der Haupteffekt}
Nach Kapitel~3:
\[ n = \frac{(1{,}96+0{,}84)^{2}\,\bigl[0{,}30\cdot0{,}70 + 0{,}20\cdot0{,}80\bigr]}{0{,}10^{2}} = \frac{7{,}85\cdot0{,}37}{0{,}01} = 291 \ \text{je Gruppe} \]

\schritt{Schritt 2 --- Der Unterschied zweier Effekte}
Er ist eine Differenz zweier Differenzen; seine Varianz ist die Summe von vier Varianzen statt zweier:
\[ n = 4\cdot291 = 1164 \ \text{je Gruppe} \]

\begin{center}
{\small
\renewcommand{\arraystretch}{1.48}
\begin{tabular}{@{}P{3.0cm}P{2.2cm}P{2.4cm}@{}}
\toprule
\textbf{Frage} & \textbf{je Gruppe} & \textbf{gesamt}\\
\midrule
Wirkt es? & $291$ & $582$\\
Wirkt es verschieden? & $1164$ & $\mathbf{4656}$\\
\bottomrule
\end{tabular}}
\end{center}

\begin{center}
\begin{minipage}{0.9\textwidth}
\textbf{Achtmal so viele Patienten für die Frage nach dem Unterschied.}
\end{minipage}
\end{center}

\begin{center}
\begin{minipage}{0.88\textwidth}
\kursiv{\textbf{Deshalb sind Untergruppenbefunde aus Studien, die für den Haupteffekt geplant wurden, fast immer unzuverlässig} --- sie haben für diese Frage eine Trennschärfe von vielleicht zwanzig Prozent. Ein \enquote{kein Effekt in dieser Untergruppe} heißt dann meist nur: zu wenige Patienten.}
\end{minipage}
\end{center}
\end{aufgabe}

\begin{aufgabe}
\boxtitel{Aufgabe 14.2 --- Behandeln oder nicht?}
Mit dem Effektverlauf aus Abschnitt 14.1: Welche Politik ist die beste, wenn die Behandlung Nebenwirkungen im Wert von $0{,}03$ Punkten hat?

\loesung
\schritt{Schritt 1 --- Die Schwelle}
Behandeln lohnt, solange $\tau(\text{Alter}) < -0{,}03$:
\[ -0{,}20 + 0{,}004\,(\text{Alter}-40) < -0{,}03 \quad\Longrightarrow\quad \text{Alter} < 82{,}5 \]

\schritt{Schritt 2 --- Der Gewinn der drei Politiken}
\begin{center}
{\small
\renewcommand{\arraystretch}{1.48}
\begin{tabular}{@{}P{2.8cm}P{2.4cm}P{2.4cm}@{}}
\toprule
\textbf{Politik} & \textbf{Nutzen} & \\
\midrule
niemanden & $0$ & \\
alle & $-0{,}100+0{,}030 = -0{,}070$ & \\
unter $82{,}5$ & $\mathbf{-0{,}072}$ & \textcolor{gruen}{beste}\\
\bottomrule
\end{tabular}}
\end{center}

\kursiv{Zur letzten Zeile:} $85$ Prozent der Patienten sind jünger als $82{,}5$; ihr mittlerer Effekt beträgt $-0{,}115$, abzüglich $0{,}03$ ergibt $0{,}85\cdot(-0{,}085) = -0{,}072$; bei den Älteren geht nichts verloren, weil sie unbehandelt bleiben.

\schritt{Schritt 3 --- Die Einordnung}
\begin{center}
\begin{minipage}{0.88\textwidth}
\kursiv{\textbf{Der Gewinn durch gezieltes Behandeln ist hier klein} --- zwei Tausendstel gegenüber \enquote{alle behandeln}. Das ist typisch: Solange der Effekt überall dasselbe Vorzeichen hat, bringt die Auswahl wenig. Lohnend wird sie erst, wenn die Behandlung bei einem Teil der Patienten \kursiv{schadet} --- und dann hängt alles daran, ob man diesen Teil zuverlässig erkennt.}
\end{minipage}
\end{center}

\schritt{Schritt 4 --- Und die Unsicherheit}
\begin{center}
\begin{minipage}{0.88\textwidth}
\kursiv{Die geschätzte Schwelle von $82{,}5$ Jahren hat selbst ein Vertrauensintervall, und zwar ein breites --- in Aufgabe 14.1 wurde sichtbar, wie teuer die Schätzung von Unterschieden ist. \textbf{Eine Politik, die an einer ungenau geschätzten Schwelle hängt, kann schlechter sein als gar keine Auswahl.} Man vergleicht deshalb Politiken nicht am geschätzten Optimum, sondern an ihrem Nutzen unter Unsicherheit.}
\end{minipage}
\end{center}
\end{aufgabe}

\begin{aufgabe}
\boxtitel{Aufgabe 14.3 --- Skala und Heterogenität}
Eine Behandlung senkt das Risiko überall um $25$ Prozent relativ. Ist der Effekt heterogen?

\loesung
\schritt{Schritt 1 --- Die absoluten Effekte}
\begin{center}
{\small
\renewcommand{\arraystretch}{1.46}
\begin{tabular}{@{}P{2.2cm}P{2.0cm}P{2.0cm}P{2.2cm}@{}}
\toprule
\textbf{Grundrisiko} & \textbf{mit} & \textbf{Differenz} & \textbf{Zahl nötig}\\
\midrule
$0{,}40$ & $0{,}30$ & $-0{,}10$ & $10$\\
$0{,}20$ & $0{,}15$ & $-0{,}05$ & $20$\\
$0{,}05$ & $0{,}0375$ & $-0{,}0125$ & $80$\\
\bottomrule
\end{tabular}}
\end{center}
\kursiv{\enquote{Zahl nötig} ist die Zahl der Behandelten je verhindertem Ereignis.}

\schritt{Schritt 2 --- Die Antwort}
\begin{center}
\begin{minipage}{0.9\textwidth}
\textbf{Auf der Verhältnisskala nein, auf der Differenzskala ja --- achtfach.}
\end{minipage}
\end{center}

\schritt{Schritt 3 --- Was für die Praxis zählt}
\begin{center}
\begin{minipage}{0.88\textwidth}
\kursiv{Die letzte Spalte entscheidet über den Einsatz: \textbf{Bei einem Grundrisiko von fünf Prozent muss man achtzig Patienten behandeln, um ein Ereignis zu verhindern} --- bei vierzig Prozent nur zehn. Ist die Behandlung teuer oder belastend, so lohnt sie im einen Fall und im anderen nicht, obwohl sie \enquote{gleich gut wirkt}.}
\end{minipage}
\end{center}

\schritt{Schritt 4 --- Die Folgerung für Leitlinien}
\begin{center}
\begin{minipage}{0.88\textwidth}
\kursiv{\textbf{Eine Leitlinie sollte den relativen Effekt aus Studien mit dem individuellen Grundrisiko aus einem Vorhersagemodell verbinden} --- statt Untergruppen nach Alter oder Geschlecht zu bilden. Damit treffen sich Band 22 und dieser Band an der Stelle, auf die es klinisch ankommt: Die Vorhersage sagt, wie gefährdet jemand ist; der kausale Schluss sagt, was die Behandlung daran ändert.}
\end{minipage}
\end{center}
\end{aufgabe}

\begin{geschichte}
\boxtitelg{Die Untergruppe der Zwillinge}
In einer großen Herzstudie der 1980er Jahre wurde die Wirkung von Aspirin nach Sternzeichen ausgewertet.

\begin{center}
\begin{minipage}{0.9\textwidth}
\kursiv{Bei Zwillingen und Waagen zeigte sich kein Nutzen, in allen anderen Zeichen ein deutlicher. \textbf{Die Verfasser druckten das Ergebnis ausdrücklich ab, um zu zeigen, wie wenig Untergruppenanalysen wert sind.} Der Scherz wird seither in jeder Methodenvorlesung erzählt --- und die Praxis der Untergruppenauswertung hat er nicht beendet.}
\end{minipage}
\end{center}

\schritt{Und ein ernster Fall}
\begin{center}
\begin{minipage}{0.9\textwidth}
\kursiv{Mehrfach wurden Medikamente aufgrund von Untergruppenbefunden gezielt eingesetzt und später widerlegt --- oder umgekehrt Patientengruppen ausgeschlossen, denen das Mittel geholfen hätte. \textbf{Die Zulassungsbehörden verlangen deshalb heute, Untergruppen vorab zu benennen} und den Wechselwirkungsterm zu prüfen, statt Gruppen getrennt auszuwerten.}
\end{minipage}
\end{center}
\end{geschichte}

\begin{forschung}
\boxtitelp{Einzelfall heute}
\begin{center}
{\small
\renewcommand{\arraystretch}{1.45}
\begin{tabular}{@{}P{3.0cm}P{5.2cm}@{}}
\toprule
\textbf{Kausale Wälder} & Heterogenität ohne vorgegebene Form, mit Vertrauensintervallen.\\
\textbf{Politiklernen} & Behandlungsregeln, die den Nutzen unter Nebenbedingungen maximieren.\\
\textbf{Risikobasierte Auswertung} & Studien nach vorhergesagtem Grundrisiko schichten.\\
\bottomrule
\end{tabular}}
\end{center}

\schritt{Zur dritten Zeile}
\begin{center}
\begin{minipage}{0.88\textwidth}
\kursiv{Man teilt die Studienteilnehmer nach ihrem \kursiv{vorhergesagten} Risiko in Viertel und berichtet den absoluten Effekt je Viertel. \textbf{Das ist eine einzige, vorab festgelegte Untergruppenbildung statt vieler} --- und sie trifft die klinisch wichtigste Quelle von Heterogenität. Mehrere Fachgesellschaften empfehlen dieses Vorgehen inzwischen als Standard.}
\end{minipage}
\end{center}
\end{forschung}

\begin{merksatz}
\boxtitel{Was Sie jetzt können}
\begin{center}
\renewcommand{\arraystretch}{1.5}
\begin{tabular}{@{}P{3.0cm}P{5.6cm}@{}}
\toprule
\textbf{$\tau(\vek z)$} & Durchschnitt über Gleichartige, nicht der Einzelne\\
\textbf{Drei Wege} & Untergruppen, Wechselwirkung, kausale Wälder\\
\textbf{Ehrlichkeit} & teilen und schätzen getrennt\\
\textbf{Achtfach} & Fallzahl für Unterschiede\\
\textbf{Skala} & Heterogenität hängt an ihr\\
\bottomrule
\end{tabular}
\end{center}

\begin{center}
\begin{minipage}{0.88\textwidth}
\kursiv{Damit endet Teil~V. Das letzte Kapitel zieht Bilanz --- und kehrt ein letztes Mal zur Tabelle aus Kapitel~1 zurück.}
\end{minipage}
\end{center}
\end{merksatz}
